\documentclass[11pt,a4paper]{article}

\usepackage[margin=1in]{geometry}
\usepackage[hidelinks]{hyperref}
\usepackage{enumitem}
\usepackage{booktabs}
\usepackage{longtable}
\usepackage{xcolor}
\usepackage{listings}

\definecolor{codebg}{HTML}{F5F5F5}
\definecolor{codecomment}{HTML}{5C6B73}
\definecolor{codekey}{HTML}{0F4C75}
\definecolor{codestr}{HTML}{B8336A}

\lstdefinestyle{shell}{
  backgroundcolor=\color{codebg},
  basicstyle=\ttfamily\footnotesize,
  commentstyle=\color{codecomment},
  keywordstyle=\color{codekey}\bfseries,
  stringstyle=\color{codestr},
  showstringspaces=false,
  breaklines=true,
  frame=single,
  framesep=4pt,
  framerule=0pt,
  xleftmargin=6pt,
  xrightmargin=6pt,
  language=bash
}
\lstset{style=shell}

\title{\textbf{MovieSentiment --- Quickstart}\\[2pt]
       \large How to Open and Run the Project from Scratch}
\author{Soumya Sarkar}
\date{\today}

\begin{document}

\maketitle

\begin{abstract}
\noindent This guide is written for somebody who has just cloned the repository and
wants to get the project running on their own machine, end to end, without prior
context. It assumes the reader is comfortable with a command line but not
necessarily with MLOps tooling. Every step lists the prerequisite, the exact
command to type, and what should happen if the step succeeds. By the end of the
guide a reader will have the API running locally on
\texttt{http://localhost:8000}, will have served a real prediction from an INT8
ONNX model, and will know where to look when something breaks.
\end{abstract}

\tableofcontents
\newpage

% ----------------------------------------------------------------------------
\section{What Are You Actually Running?}
% ----------------------------------------------------------------------------

MovieSentiment classifies IMDb-style movie reviews as positive or negative. The
fast path through the project is:

\begin{enumerate}[leftmargin=*]
\item Install Python and the project's dependencies.
\item Pull the pre-trained ONNX model files (so you do not need a GPU).
\item Start the FastAPI service.
\item Send it some text.
\end{enumerate}

\noindent That is the five-minute version. The remainder of this document walks
through each step slowly, then covers Docker, running the full training pipeline
yourself, executing tests, and deploying to AWS.

% ----------------------------------------------------------------------------
\section{Prerequisites}
% ----------------------------------------------------------------------------

\subsection{Required Software}

\begin{longtable}{@{}lp{0.30\linewidth}p{0.4\linewidth}@{}}
\toprule
\textbf{Tool} & \textbf{Minimum version} & \textbf{Why you need it} \\ \midrule \endhead
Python   & 3.10 (3.11 recommended)        & The language the project is written in. \\
Git      & any modern version             & To clone the repository. \\
pip      & ships with Python              & Bootstraps everything else. \\
uv       & latest                         & A fast pip replacement; installed with pip. \\
Docker   & 20.10 or later (optional)      & For the local Prometheus + Grafana stack. \\
DVC      & 3.48 or later (optional)       & For pulling the pre-trained model files
                                            from S3. \\
AWS CLI  & latest (optional)              & Only needed if you plan to deploy or pull
                                            DVC artefacts from the project's S3 remote. \\
\bottomrule
\end{longtable}

\subsection{How to Check What You Have}

\begin{lstlisting}
python --version
git --version
docker --version            # only if you plan to use the docker-compose stack
\end{lstlisting}

\noindent If \texttt{python} maps to Python 2 on your machine, use \texttt{python3}
in every command below.

\subsection{Installing the Missing Pieces}

\begin{itemize}[leftmargin=*]
\item \textbf{Python}: download from \texttt{https://www.python.org/downloads/} (Windows,
      macOS, Linux) or use your distribution's package manager.
\item \textbf{Git}: \texttt{https://git-scm.com/downloads}.
\item \textbf{Docker Desktop} (optional): \texttt{https://www.docker.com/products/docker-desktop/}.
\end{itemize}

% ----------------------------------------------------------------------------
\section{Step 1 --- Clone the Repository}
% ----------------------------------------------------------------------------

\begin{lstlisting}
git clone https://github.com/Cryptic2-0/Project moviesentiment
cd moviesentiment
\end{lstlisting}

\noindent The working tree now looks like this:

\begin{lstlisting}
moviesentiment/
  src/moviesentiment/      # the installable package
  tests/                   # pytest suite
  deploy/                  # Dockerfile, ECS task, Lambda handler, compose
  scripts/                 # SageMaker, Locust, lock-file scripts
  frontend/                # the single-page demo UI
  docs/                    # this file, model card, datasheet, benchmarks
  dvc.yaml                 # the pipeline graph
  params.yaml              # hyperparameters
  pyproject.toml           # dependencies
  Makefile                 # day-to-day shortcuts
\end{lstlisting}

% ----------------------------------------------------------------------------
\section{Step 2 --- Create a Virtual Environment}
% ----------------------------------------------------------------------------

The virtual environment isolates the project's Python packages from your system
Python, so nothing you install for this project leaks into other work.

\subsection{macOS / Linux}

\begin{lstlisting}
python -m venv .venv
source .venv/bin/activate
\end{lstlisting}

\subsection{Windows (PowerShell)}

\begin{lstlisting}
python -m venv .venv
.\.venv\Scripts\Activate.ps1
\end{lstlisting}

\noindent After this command your shell prompt should be prefixed with
\texttt{(.venv)}.

% ----------------------------------------------------------------------------
\section{Step 3 --- Install Dependencies}
% ----------------------------------------------------------------------------

The project uses \texttt{uv} for installs because it is dramatically faster than
plain \texttt{pip} on a large dependency tree. Bootstrap it once:

\begin{lstlisting}
pip install uv
\end{lstlisting}

\noindent Then install the project in editable mode with the development extras:

\begin{lstlisting}
uv pip install -e ".[dev]"
\end{lstlisting}

\noindent What ``editable mode'' means: edits you make under \texttt{src/} take
effect immediately, without re-installing. The \texttt{[dev]} extra adds
ruff, black, mypy, pytest, hypothesis, locust, pre-commit, bandit, and pip-audit.

\subsection{Optional Extras}

\begin{longtable}{@{}lp{0.6\linewidth}@{}}
\toprule
\textbf{Extra} & \textbf{What it adds} \\ \midrule \endhead
\texttt{train}   & SageMaker SDK, boto3, peft, accelerate. Needed only if you
                   plan to fine-tune the transformer yourself. \\
\texttt{explain} & captum, for gradient-based interpretability experiments. \\
\texttt{monitor} & nannyml, great-expectations, for production performance
                   estimation. \\
\bottomrule
\end{longtable}

\noindent Install any extra with:

\begin{lstlisting}
uv pip install -e ".[dev,train]"
\end{lstlisting}

% ----------------------------------------------------------------------------
\section{Step 4 --- Obtain the Model Files}
% ----------------------------------------------------------------------------

The repository tracks the ONNX model artefacts with DVC, not Git. You have three
ways to get them:

\subsection{Option A --- Pull from the Project's S3 Remote (Recommended)}

\textbf{Prerequisite:} AWS credentials with read access to the DVC bucket. If you
were given keys, set them in your shell first:

\begin{lstlisting}
export AWS_ACCESS_KEY_ID=AKIA...
export AWS_SECRET_ACCESS_KEY=...
# Windows PowerShell:
# $env:AWS_ACCESS_KEY_ID = "AKIA..."
# $env:AWS_SECRET_ACCESS_KEY = "..."

dvc pull models/distilbert_onnx_int8.dvc models/distilbert_onnx.dvc
\end{lstlisting}

\noindent After the pull you should see
\texttt{models/distilbert\_onnx\_int8/model.onnx} (about 64\,MB) and
\texttt{models/distilbert\_onnx/model.onnx} (about 256\,MB).

\subsection{Option B --- Export the ONNX Files Yourself}

If you have no S3 access but do have time, let the project download a base
DistilBERT from HuggingFace and export it locally:

\begin{lstlisting}
python -m moviesentiment.models.onnx_export
# or: make export-onnx
\end{lstlisting}

\noindent The script will print \texttt{models/distilbert/ not found --- exporting
base distilbert-base-uncased} and produce both FP32 and INT8 versions, plus the
benchmark report at \texttt{docs/benchmarks.md}. The base model is \emph{not}
fine-tuned, so accuracy on actual reviews will be poor; this option is for
verifying that the serving path works.

\subsection{Option C --- Train the Model from Scratch}

See the section ``Running the Training Pipeline'' further down. This requires
either a GPU on your laptop or a SageMaker account.

% ----------------------------------------------------------------------------
\section{Step 5 --- Start the API Locally}
% ----------------------------------------------------------------------------

With the model in place, start the FastAPI service:

\begin{lstlisting}
make serve
# equivalent to:
# uvicorn moviesentiment.serve.api:app --host 0.0.0.0 --port 8000 --reload
\end{lstlisting}

\noindent You should see structured-JSON log lines, including one with
\texttt{"event": "model\_loaded"}. The service is listening on port 8000.

\subsection{Useful URLs}

\begin{longtable}{@{}lp{0.6\linewidth}@{}}
\toprule
\textbf{URL} & \textbf{What it shows} \\ \midrule \endhead
\texttt{http://localhost:8000/docs}        & Auto-generated Swagger UI -- try every endpoint here. \\
\texttt{http://localhost:8000/healthz}     & Always returns \texttt{\{"status":"ok"\}} once running. \\
\texttt{http://localhost:8000/readyz}      & 200 once the model is loaded, 503 otherwise. \\
\texttt{http://localhost:8000/version}     & Model name, stage, and the git SHA. \\
\texttt{http://localhost:8000/metrics}     & Prometheus exposition. \\
\texttt{http://localhost:8000/ui/}         & The static demo frontend (if bundled). \\
\bottomrule
\end{longtable}

% ----------------------------------------------------------------------------
\section{Step 6 --- Send Your First Prediction}
% ----------------------------------------------------------------------------

In a second terminal, with the server still running:

\begin{lstlisting}
curl -X POST http://localhost:8000/predict \
     -H "Content-Type: application/json" \
     -d '{"texts": ["A complete masterpiece.",
                    "Two hours I will never get back."]}'
\end{lstlisting}

\noindent The response will look like:

\begin{lstlisting}
{
  "predictions": [
    {"text": "A complete masterpiece.",
     "label": "positive",
     "confidence": 0.987},
    {"text": "Two hours I will never get back.",
     "label": "negative",
     "confidence": 0.962}
  ]
}
\end{lstlisting}

\subsection{Try the Explainer}

\begin{lstlisting}
curl -X POST http://localhost:8000/explain \
     -H "Content-Type: application/json" \
     -d '{"text":"The film was breathtakingly beautiful.","top_k":5}'
\end{lstlisting}

\noindent You will receive the top-5 tokens ranked by how much they supported the
predicted label, computed by dropping each token in turn and measuring the change
in confidence.

\subsection{Limits to Be Aware Of}

\begin{itemize}[leftmargin=*]
\item \texttt{POST /predict} accepts at most 32 texts per request.
\item Each text must be at most 5000 characters long.
\item \texttt{POST /predict} is rate-limited to 60 requests per minute per IP.
\item \texttt{POST /explain} is rate-limited to 10 requests per minute per IP.
\end{itemize}

% ----------------------------------------------------------------------------
\section{Configuration via Environment Variables}
% ----------------------------------------------------------------------------

Every tuneable setting is read by the Pydantic \texttt{Settings} class in
\texttt{src/moviesentiment/config.py}. The variables all use the prefix
\texttt{MS\_}; you can either export them in your shell or write them into a
\texttt{.env} file at the repository root.

\begin{longtable}{@{}lp{0.32\linewidth}p{0.34\linewidth}@{}}
\toprule
\textbf{Variable} & \textbf{Default} & \textbf{Effect} \\ \midrule \endhead
\texttt{MS\_MAX\_TEXT\_LENGTH}      & 5000          & Per-text character cap. \\
\texttt{MS\_MAX\_BATCH\_SIZE}       & 32            & Per-request batch cap. \\
\texttt{MS\_CORS\_ALLOW\_ORIGINS}   & \texttt{*}    & Comma-separated CORS allowlist;
                                                     an explicit list auto-enables
                                                     credentials. \\
\texttt{MS\_HF\_REVISION}           & empty         & Pin a HuggingFace commit SHA;
                                                     empty means \texttt{main}
                                                     (dev only). \\
\texttt{MS\_MLFLOW\_TRACKING\_URI}  & \texttt{sqlite:///mlflow.db} & Where MLflow
                                                                    writes runs. \\
\texttt{MS\_MODEL\_NAME}            & \texttt{moviesentiment-classifier} & Registry name. \\
\texttt{MS\_MODEL\_STAGE}           & \texttt{Production} & Which stage to load. \\
\texttt{OTEL\_EXPORTER\_OTLP\_ENDPOINT} & unset    & Set to enable distributed tracing. \\
\texttt{MS\_SQS\_QUEUE\_URL}, \texttt{MS\_JOBS\_TABLE} & unset & Required for the
                                                                async path; without
                                                                them \texttt{/predict/async}
                                                                returns 503. \\
\bottomrule
\end{longtable}

% ----------------------------------------------------------------------------
\section{Running the Local Observability Stack}
% ----------------------------------------------------------------------------

If you have Docker installed, \texttt{deploy/docker-compose.yml} brings up the API
plus Prometheus and Grafana with the dashboard pre-provisioned.

\begin{lstlisting}
docker compose -f deploy/docker-compose.yml up --build
# or:
make docker-up
\end{lstlisting}

\noindent After the stack is up:

\begin{longtable}{@{}lp{0.5\linewidth}@{}}
\toprule
\textbf{Service} & \textbf{URL} \\ \midrule \endhead
API        & \texttt{http://localhost:8000} \\
Prometheus & \texttt{http://localhost:9090} \\
Grafana    & \texttt{http://localhost:3000} (anonymous Admin auto-enabled) \\
\bottomrule
\end{longtable}

\noindent The Grafana dashboard reads the four custom metrics
(\texttt{prediction\_confidence}, \texttt{prediction\_class\_total},
\texttt{model\_version\_info}, \texttt{model\_estimated\_f1}) plus the default
\texttt{prometheus-fastapi-instrumentator} histograms.

% ----------------------------------------------------------------------------
\section{Running the Training Pipeline}
% ----------------------------------------------------------------------------

The full pipeline -- scrape, clean, validate, split, train -- is declared in
\texttt{dvc.yaml}. To execute everything:

\begin{lstlisting}
dvc repro
\end{lstlisting}

\noindent DVC walks the DAG and only re-runs stages whose dependencies changed.
The first run is the expensive one (downloading 50{,}000 IMDb reviews and fine-tuning
DistilBERT for three epochs). On a laptop without a GPU the transformer stage will
take hours; the baseline stage finishes in minutes.

\subsection{Running Individual Stages}

\begin{lstlisting}
python -m moviesentiment.cli scrape
python -m moviesentiment.cli clean
python -m moviesentiment.data.validate \
       data/interim/clean.parquet metrics/data_quality.json
python -m moviesentiment.cli split
python -m moviesentiment.cli train baseline
python -m moviesentiment.cli train transformer
\end{lstlisting}

\subsection{Inspecting Results in MLflow}

\begin{lstlisting}
mlflow ui --backend-store-uri sqlite:///mlflow.db
# open http://localhost:5000
\end{lstlisting}

\subsection{Rendering Metric Tables in the Terminal}

\begin{lstlisting}
python -m moviesentiment.cli metrics
\end{lstlisting}

\noindent Walks \texttt{metrics/*.json} and prints a Rich-formatted table.

\subsection{Drift Reports}

\begin{lstlisting}
python -m moviesentiment.cli drift
# HTML report written to docs/drift_reports/<today>.html
\end{lstlisting}

% ----------------------------------------------------------------------------
\section{Running the Tests}
% ----------------------------------------------------------------------------

\begin{lstlisting}
make test          # fast tests, coverage gate at 85%
make test-slow     # ONNX export e2e tests (minutes long)
\end{lstlisting}

\noindent Behind the scenes \texttt{make test} runs
\texttt{pytest -m "not slow" --cov=moviesentiment --cov-fail-under=85 -v}. A failing
coverage gate is a real failure -- it will turn CI red.

\subsection{Linting and Type Checking}

\begin{lstlisting}
make lint           # ruff + black
make fmt            # auto-fix the same
make type-check     # mypy on src and tests
make pre-commit     # the full pre-commit suite
make ci-local       # the bundle CI runs: lint + type + test
\end{lstlisting}

\subsection{The Locust Load Test}

\begin{lstlisting}
# Terminal 1: keep the API running.
make serve

# Terminal 2: hammer it for two minutes at 50 concurrent users.
make loadtest
# Report at docs/load_test_report.html
\end{lstlisting}

% ----------------------------------------------------------------------------
\section{Deploying to AWS (Reference)}
% ----------------------------------------------------------------------------

This is not required to run the project. It is here so a reader who wants to put
the service on the public internet knows what the project actually does in CI.

\subsection{One-Time Setup}

\begin{enumerate}[leftmargin=*]
\item Create an ECR repository in \texttt{ap-southeast-2} named
      \texttt{moviesentiment}.
\item Create an ECS cluster called \texttt{moviesentiment} and a service of the
      same name backed by the task definition in
      \texttt{deploy/ecs-task-definition.json} (Fargate, ARM64, 0.25 vCPU,
      1024 MB).
\item Create an S3 bucket for the DVC remote and grant the IAM user / role
      read+write.
\item Store \texttt{AWS\_ACCESS\_KEY\_ID}, \texttt{AWS\_SECRET\_ACCESS\_KEY} as
      GitHub Actions secrets so \texttt{ci.yml} can pull DVC artefacts and push to
      ECR.
\end{enumerate}

\subsection{Pushing a New Build}

Each push to \texttt{main} triggers \texttt{.github/workflows/ci.yml}, which lints,
tests, builds the ARM64 image with \texttt{GIT\_SHA=\$\{\{ github.sha \}\}}, pushes
to ECR plus GHCR, and then calls:

\begin{lstlisting}
aws ecs update-service \
    --cluster moviesentiment \
    --service moviesentiment \
    --force-new-deployment \
    --region ap-southeast-2
\end{lstlisting}

\subsection{Smoke-Testing the Live Endpoint}

\begin{lstlisting}
make smoke-test
# Resolves the running task's public IP, then curls /predict.
\end{lstlisting}

\subsection{Monthly Cost}

About \$6/mo for the Fargate task (Graviton ARM, lowest tier). DynamoDB, SQS, and
Lambda for the async path are inside their free tiers under portfolio traffic.

% ----------------------------------------------------------------------------
\section{Common Issues}
% ----------------------------------------------------------------------------

\begin{longtable}{@{}p{0.42\linewidth}p{0.5\linewidth}@{}}
\toprule
\textbf{Symptom} & \textbf{What it usually means} \\ \midrule \endhead
\texttt{/readyz} returns 503 forever. & The ONNX file is not in
                                        \texttt{models/distilbert\_onnx\_int8/}.
                                        Pull via DVC or run \texttt{make export-onnx}. \\
\texttt{/predict} returns 422. & Either the batch exceeds 32 texts or one of the
                                 texts exceeds 5000 characters. \\
\texttt{429 Too Many Requests}. & The slowapi rate limit fired. Wait a minute. \\
\texttt{mlflow ui} shows no runs. & Make sure \texttt{MS\_MLFLOW\_TRACKING\_URI} is
                                    \texttt{sqlite:///mlflow.db} and that you ran
                                    training in the same working directory. \\
\texttt{dvc pull} fails with a 403. & AWS credentials are missing or do not have
                                      read on the bucket. \\
Docker compose hangs on the API. & First build downloads the model layer; expect
                                   a one-time slow build. \\
\texttt{pytest} reports coverage failure. & You have probably touched a serving
                                            module without updating its tests --
                                            coverage gate is 85\%. \\
\bottomrule
\end{longtable}

% ----------------------------------------------------------------------------
\section{Where to Look Next}
% ----------------------------------------------------------------------------

\begin{itemize}[leftmargin=*]
\item \texttt{workflow.tex} -- the companion document; explains \emph{what} the
      project does end to end and why each component exists.
\item \texttt{docs/model\_card.md} -- intended use, per-slice F1, fairness notes,
      environmental impact estimate.
\item \texttt{docs/datasheet.md} -- the IMDb 50K dataset provenance and bias notes.
\item \texttt{docs/benchmarks.md} -- raw FP32 vs INT8 latency numbers.
\item \texttt{docs/loadtest.md} -- Locust reference numbers, bottleneck breakdown.
\item \texttt{docs/grafana\_cloud\_setup.md} -- one-time OTLP export setup.
\item \texttt{docs/future\_improvements.md} -- deferred items and their pickup plans.
\item \texttt{SECURITY.md} -- threat model, mitigations, accepted residual risks.
\end{itemize}

\end{document}
